% !TEX root = ../main.tex
\lecture{2}{2026-09-21}{Woche 2}

\section{Investitionsrechnung}
\textbf{Definition} \textit{Investition} Aktiventausch vom flüssigen Mittel (Geld) in Anlagevermögen [Kapital]. \textit{Siehe figures Bilanz\_Investition.png}
\textbf{Zentrale Fragen der Investitionsrechnung} Wie viel und worin soll investiert werden?
\subsection{Rechenverfahren}
\subsubsection{Statische Investitionsrechenverfahren}

\begin{itemize}
    \item Basisdaten
    \item Kaufpreis (I)
    \item Liquidationswert (L)
    \subitem Endwert Devaluation nach volle Nutzungsdauer
    \item Kapitalzinsen
    \item Nutzungsdauer (n)
    \item Jaehrliche Abschreibung (Devaluation) = I/n 
    \item Kapitalzinsen / Jahr
\end{itemize}

\textbf{Defnition} \textit{Kalkulationszins} ist Reditemöglichkeit einer Alternativeinvestition.
\begin{theorem}
    Kalkulatorische Zinsen (Betrag, kein Cash Flow) = Gebundenes Kapital (Avg Kapital dass alternativ arrenditiert) * Kapitalzinssatz 
\end{theorem}
\break \textbf{Beispiel} \break 
Bei einer Investition von 240, die über 8 Jahre ihren ganzen Wert verliert, und Kapital/Kalkulationszins von 10\%.
Sind mit durchschnittliche 10\% Zinsen auf (240-0)/2 oder 120 pro Jahr zu rechnen. Diese sind entgangene Rendite (Opportunitätskosten der Investition)

\begin{itemize}
    \item Kostenrechnung (Summe von)
        \subitem Betriebskosten
        \subitem Lohnkosten
        \subitem Abschreibung (Devaluation)
        \subitem Kalkulationszins
    \item Gewinnsrechnung
        \subitem Umsatz - Kosten 
    \item Rentabilitätsrechnung (Summe von)
        \subitem Gewinn
        \subitem Kalkulationszins
        \subitem 
    \item Amortisationsrechnung
        \subitem Zusätzliche Einnahmen - Ausgaben pro jahr
        \subitem überschüssige Einnahmen pro jahr
        \subitem Amortisationszeit = Investition / ü. Einnahmen = X Jahre
\end{itemize}

\textbf{Zusammenfassung statische}
\begin{itemize}
    \item engen Bezug Rechnungswesen
\end{itemize}